\chapter{Introduction}

\section{Industrial Context}

Thermal process control is ubiquitous in modern industry. Whether in chemical reactors,
food processing lines, pharmaceutical production, or power generation systems, the ability
to regulate temperatures precisely and reliably directly determines product quality, energy
efficiency, and process safety. A plate heat exchanger, for example, appears at every scale
of industrial production: from pasteurisation of beverages to cooling of power electronics
and conditioning of building HVAC systems.

The supervision of such processes has historically relied on dedicated SCADA
(Supervisory Control and Data Acquisition) systems — proprietary software suites
connected to programmable logic controllers via industrial fieldbus protocols. While reliable
and certified, classical SCADA systems carry significant costs in licensing, hardware, and
engineering time. Their interfaces, designed for control-room workstations, are not naturally
suited to the web-connected, multi-device environments that modern facilities increasingly
demand.

Over the past decade, the convergence of industrial automation and information technology
has given rise to a new paradigm often called Industrial~IoT (IIoT) or Industry~4.0.
In this paradigm, PLCs and sensors publish data to web-accessible services; dashboards run
in standard browsers; and control logic can be deployed as containerised software on hardware
directly integrated into the control cabinet or operator panel. Siemens' Industrial Edge
platform is a concrete realisation of this vision: it allows Docker containers to run natively
on a SIMATIC Unified panel, inches from the process they supervise.

This project is situated at this intersection. The installation used is a laboratory-scale
thermal rig at ICAM Engineering School, equipped with a Siemens S7-1500 PLC, a plate heat
exchanger, an electric furnace, and a SIMATIC Unified 7\textquotedblright{} touch panel.
It provides a representative testbed for exploring how modern web technologies---Python,
Flask, Socket.IO, Chart.js, Docker---can be combined with an industrial-grade controller
to deliver a complete supervision and control solution without the cost or rigidity of a
classical SCADA system.

\section{Problem Statement}

The installation under study is a plate heat exchanger in which a hot fluid (circuit~F1),
pre-heated by an electric furnace, transfers heat to a cold counter-current fluid
(circuit~F2). The primary control objective is to maintain the cold outlet temperature
$T_{out2,HE1}$ at a desired setpoint, despite disturbances arising from variations in
furnace temperature or cold circuit flow rate.

A straightforward approach would be to apply a single PID controller closing the loop
directly around $T_{out2,HE1}$ by manipulating the hot flow setpoint $F_{1,SP}$.
However, this approach faces a fundamental limitation: the dynamic path from $F_{1,SP}$
to $T_{out2,HE1}$ passes through two sequential processes --- the furnace
($F_{1,SP} \rightarrow T_{in1,HE1}$) and the heat exchanger
($T_{in1,HE1} \rightarrow T_{out2,HE1}$). The combined transfer function exhibits a
relatively large apparent dead time and a long settling time. A controller tuned
aggressively enough to achieve fast setpoint tracking will tend to oscillate; one tuned
conservatively will be too slow to reject disturbances before they affect $T_{out2}$.

\textit{Cascade control} offers an elegant solution to this problem. By introducing an
intermediate measurement~--- the heat exchanger hot inlet temperature $T_{in1,HE1}$~---
and dedicating a fast inner control loop to regulate it, disturbances in the hot circuit
are corrected before they propagate to the cold outlet. An outer loop then adjusts the
$T_{in1}$ setpoint to guide $T_{out2}$ toward the desired value. This structure improves
both setpoint tracking and disturbance rejection compared to a single-loop strategy.

A second problem addressed is the absence of any web-based supervision interface for this
installation. Operators currently have no real-time visibility into the process state from
a network-connected device. The panel displays a static WinCC Unified screen, but no data
logging, trend charts, or remote accessibility exist. Building such an interface, deployable
both on a PC and directly on the panel via Industrial Edge, represents a significant
operational improvement.

\section{Objectives}

This project pursues three main objectives:

\begin{enumerate}[leftmargin=*, itemsep=4pt]

  \item \textbf{Real-time web supervision.}
    Design and implement a web application that communicates with the S7-1500 PLC in real
    time, displays all relevant process variables (temperatures, flow rates, valve positions,
    furnace state), provides setpoint controls, and logs data to a persistent file. The
    interface must be accessible both from a workstation and from the embedded browser of
    the SIMATIC Unified 7\textquotedblright{} panel.

  \item \textbf{Experimental process identification.}
    Conduct an open-loop step-response test on the inner process ($F_{1,SP} \rightarrow
    T_{in1,HE1}$) to extract the parameters of a First-Order Plus Time Delay (FOPTD)
    model. Use these parameters to calculate PI controller gains according to Internal Model
    Control (IMC) tuning rules, providing three levels of aggressiveness for the engineer.

  \item \textbf{Cascade control implementation.}
    Implement a cascade controller in Python comprising: (a)~an inner PI loop with
    proportional action on measurement (PI+P), regulating $T_{in1,HE1}$ by manipulating
    $F_{1,SP}$; and (b)~an outer proportional controller, converting a $T_{out2,HE1}$
    setpoint into a $T_{in1,HE1}$ setpoint. The implementation must include bumpless mode
    transfer, anti-windup, and a three-mode operating interface (Manual / Auto-Inner /
    Auto-Full).

\end{enumerate}

\section{Thesis Structure}

\textbf{Chapter~2} describes the physical installation: the electric furnace, the plate
heat exchanger, instrumentation, and the control objectives derived from the process physics.

\textbf{Chapter~3} covers the Siemens S7-1500 PLC, its memory organisation and scan cycle,
TIA Portal V18 configuration, WinCC Unified on the 7\textquotedblright{} panel, the Open
Pipe communication mechanism, and Industrial Edge deployment.

\textbf{Chapter~4} presents the software architecture: technology choices, the multi-threaded
Flask application, the two PLC communication layers (Snap7 and Open Pipe), real-time
data flow via Socket.IO, the REST API, and the four web interface pages.

\textbf{Chapter~5} develops the theory of FOPTD process identification and IMC tuning,
then presents the experimental protocol and results.

\textbf{Chapter~6} details the cascade control design: inner PI+P controller (velocity
form, bumpless transfer, anti-windup), outer proportional controller, mode state machine,
and experimental validation.

\textbf{Chapter~7} describes the supervision interface design: the panel-optimised P\&ID
synoptic and the cascade commissioning page.

\textbf{Chapter~8} concludes with a summary of contributions, assessment of results,
current limitations, and perspectives.

\section{Scope and Limitations}

This project focuses on the HE-001 plate heat exchanger installation at ICAM. The
identification and closed-loop control work is limited to the inner process
($F_{1,SP} \rightarrow T_{in1,HE1}$). The outer loop uses a proportional controller
whose gain is set manually; full outer loop identification is left for future work.

The web application is implemented as a proof-of-concept demonstrator with an emphasis
on functionality and real-time performance. It does not include industrial certification
(IEC~61508/61511), formal safety interlocks, or redundancy mechanisms. The Open Pipe
communication protocol is validated against the available Siemens documentation; full
end-to-end validation with the WinCC Unified runtime on the panel is pending Industrial
Edge licence acquisition.
